% Shared preamble for the literature/survey LaTeX documents.
% Used as the very first line of each document:  % Shared preamble for the literature/survey LaTeX documents.
% Used as the very first line of each document:  % Shared preamble for the literature/survey LaTeX documents.
% Used as the very first line of each document:  \input{preamble}
% (\input is a TeX primitive, so it is legal before \begin{document}.)
\documentclass[11pt]{article}
\usepackage[T1]{fontenc}
\usepackage{lmodern}
\usepackage[letterpaper,margin=1in]{geometry}
\usepackage{amsmath,amssymb}
\usepackage{array,booktabs}
\usepackage{enumitem}
\usepackage[most]{tcolorbox}
\usepackage{xcolor}
\usepackage[colorlinks=true,linkcolor=teal!60!black,urlcolor=teal!60!black,citecolor=teal!60!black]{hyperref}
\usepackage[backend=biber,style=numeric-comp,sorting=ynt,maxbibnames=8,maxcitenames=2,giveninits=true]{biblatex}
\addbibresource{references.bib}
\newtcolorbox{callout}{colback=black!3,colframe=black!25,boxrule=0.4pt,arc=2pt,
  left=6pt,right=6pt,top=4pt,bottom=4pt,breakable}
\setlist{topsep=3pt,itemsep=1pt,parsep=1pt,leftmargin=1.4em}
\setlength{\parindent}{0pt}\setlength{\parskip}{0.5em}
\sloppy

% (\input is a TeX primitive, so it is legal before \begin{document}.)
\documentclass[11pt]{article}
\usepackage[T1]{fontenc}
\usepackage{lmodern}
\usepackage[letterpaper,margin=1in]{geometry}
\usepackage{amsmath,amssymb}
\usepackage{array,booktabs}
\usepackage{enumitem}
\usepackage[most]{tcolorbox}
\usepackage{xcolor}
\usepackage[colorlinks=true,linkcolor=teal!60!black,urlcolor=teal!60!black,citecolor=teal!60!black]{hyperref}
\usepackage[backend=biber,style=numeric-comp,sorting=ynt,maxbibnames=8,maxcitenames=2,giveninits=true]{biblatex}
\addbibresource{references.bib}
\newtcolorbox{callout}{colback=black!3,colframe=black!25,boxrule=0.4pt,arc=2pt,
  left=6pt,right=6pt,top=4pt,bottom=4pt,breakable}
\setlist{topsep=3pt,itemsep=1pt,parsep=1pt,leftmargin=1.4em}
\setlength{\parindent}{0pt}\setlength{\parskip}{0.5em}
\sloppy

% (\input is a TeX primitive, so it is legal before \begin{document}.)
\documentclass[11pt]{article}
\usepackage[T1]{fontenc}
\usepackage{lmodern}
\usepackage[letterpaper,margin=1in]{geometry}
\usepackage{amsmath,amssymb}
\usepackage{array,booktabs}
\usepackage{enumitem}
\usepackage[most]{tcolorbox}
\usepackage{xcolor}
\usepackage[colorlinks=true,linkcolor=teal!60!black,urlcolor=teal!60!black,citecolor=teal!60!black]{hyperref}
\usepackage[backend=biber,style=numeric-comp,sorting=ynt,maxbibnames=8,maxcitenames=2,giveninits=true]{biblatex}
\addbibresource{references.bib}
\newtcolorbox{callout}{colback=black!3,colframe=black!25,boxrule=0.4pt,arc=2pt,
  left=6pt,right=6pt,top=4pt,bottom=4pt,breakable}
\setlist{topsep=3pt,itemsep=1pt,parsep=1pt,leftmargin=1.4em}
\setlength{\parindent}{0pt}\setlength{\parskip}{0.5em}
\sloppy

\begin{document}

\section*{Group 6 --- Epistemics and vision}

\begin{callout}
This group is the one that steps back from technique and asks the meta-questions: \emph{what does it mean for a mechanistic interpretation to be correct, and where is the whole program trying to go?} Where Groups 1--5 produce machinery --- circuits, polytopes, dictionaries, attribution graphs, metrics --- Group 6 supplies the program's self-image: its vocabulary, its ambition, and the one sharp counterexample that disciplines both. There is less mathematics here than elsewhere, but it is not absent: the unfaithfulness toy model contains the single most pointed formal idea about \emph{what it means to match a mechanism rather than a function}, and we treat it precisely.
\end{callout}


\subsection*{The shared question}

The three papers form an arc --- \textbf{vocabulary}, \textbf{ambition}, \textbf{epistemic check} --- around one question: when an interpreter says "this circuit computes $X$", what would make that statement \emph{true}, and how would we know?

\textcite{mech-interp-variables-bases} supplies the conceptual scaffolding. It argues, by sustained analogy to reverse-engineering a compiled binary, that interpretability's central task is to \emph{name the variables}: to decompose high-dimensional activations into independently understandable pieces, just as a reverse engineer must segment raw memory into variables before any line of code becomes meaningful. \textcite{interpretability-dreams} supplies the ambition: a frankly speculative manifesto for what becomes possible \emph{once superposition is resolved} --- a "complete set of features" to enumerate over, larger-scale structure built on top of circuits, universality across models, and an \textbf{epistemic foundation} in which, "for small enough circuits, statements about their behaviour become questions of mathematical reasoning." \textcite{toy-model-of-mechanistic-unfaithfulness} supplies the cold shower: a fully worked example in which a sparse, monosemantic, low-error model of a computation nonetheless implements the \emph{wrong algorithm} --- the same function, a different mechanism --- and therefore generalises differently off-distribution. It also offers a partial remedy with real mathematical content (Jacobian matching).

The thesis tying them together: \textbf{mechanistic interpretability is staked on a notion of correctness --- "mechanistic faithfulness" --- that is more demanding than input--output agreement and is, as yet, largely unformalised.} The vocabulary paper says \emph{what} we are trying to recover (interpretable variables); the dreams paper says \emph{why} (because circuits could become a rigorous epistemic substrate for everything above them); the unfaithfulness paper says \emph{where this can silently fail} and gestures at the metric that would catch it. This group is, in effect, the field reflecting on its own epistemology, and it is exactly the seam where a mathematician's instinct for "what is the definition, and what is the theorem" is most useful and most missing.


\subsection*{The papers}

\subsubsection*{The program's vocabulary: \texorpdfstring{\textcite{mech-interp-variables-bases}}{Variables and bases}}

This essay establishes the metaphor the entire lineage runs on. A trained network is "in some sense a binary computer program which runs on one of the exotic virtual machines we call a neural network architecture", and so the table of correspondences is: \textbf{network parameters $\approx$ program binary}; \textbf{architecture $\approx$ the VM/processor}; \textbf{layer activations $\approx$ program state / memory}; and --- the load-bearing entry --- \textbf{a neuron or feature direction $\approx$ a variable / memory location}. Mechanistic interpretability is then reverse-engineering, and the central act is \emph{assigning meaning to positions in memory}: segmenting the activation vector into pieces that can be understood independently.

Two arguments carry the weight. First, the \textbf{curse of dimensionality} as an obstacle \emph{to interpretation}, not just to learning. A function on an exponentially large input space cannot be understood by tabulating its behaviour. But a reverse engineer escapes this because the \emph{code} is a non-exponential, finite description of behaviour; the parameters of a network are likewise a finite description, if only one could read them. (The essay anticipates and rejects the "just restrict to the data manifold" escape: the manifold of natural images or text is itself very high-dimensional --- an object can enter a scene from any direction --- so it offers no real reduction.) The interpretability program is thus the bet that the finite parameter-description can be made human-legible, decomposed into independently-reasoned parts the way memory decomposes into variables.

Second, the \textbf{privileged basis} as the thing that makes this hope tractable, and \textbf{polysemanticity} as the obstruction to it. In principle features could sit at arbitrary directions in activation space; in practice, a pointwise activation function "makes these directions natural and useful" and tends to align features with individual neurons --- exactly the basis-privileging mechanism formalised in the \emph{Landscape} survey I.4 and the \emph{ML~Concepts} primer \S1.5. This is why neuron-level interpretation is even conceivable. But many neurons are stubbornly \textbf{polysemantic}: they fire for several unrelated concepts, because the model is representing more features than it has neurons --- the symptom of superposition. So the essay names the disease (polysemanticity), locates its cause (superposition), and identifies the one architectural gift (the privileged basis from the nonlinearity) that makes any of it readable. It does not solve anything; it sets the agenda the rest of the corpus pursues.


\subsubsection*{The aspiration: \texorpdfstring{\textcite{interpretability-dreams}}{Interpretability Dreams}}

This is the manifesto, and it announces itself as "almost definitionally extremely speculative" --- a deliberate surfacing of motivating visions, not a results paper. Its spine is the claim that mechanistic interpretability is a \textbf{microscopic} science deliberately built bottom-up to be an \emph{epistemic foundation}. The argument for going bottom-up is an argument about \emph{being misled}: a top-down probe for a "dog-head detector" will succeed even in a layer that has no such feature, because upstream snout/fur/eye detectors already predict dog-heads well enough to fool a linear probe. The corrective is rigour at the bottom: "for small enough circuits, statements about their behaviour become questions of mathematical reasoning." The cost of that rigour is scope --- circuit-level statements are small --- but the wager is that statements about model behaviour can be \emph{reduced} to statements about circuits, so circuits act as a foundation one can fall back to in order to \emph{check} any higher-level claim. The essay is careful that this is a foundation in the logician's sense, not a working register: as mathematicians do not reason from the axioms day-to-day and programmers do not think in assembly, the value of the foundation is that disputes can in principle be reduced to it.

On that foundation the essay stacks its aspirations, all gated on \emph{first resolving superposition}. (i) \textbf{Feature families and equivariance}: features come organised into families parameterised by a variable (orientation, hue, scale; "token $X$ in language $Y$"), and crucially the \emph{circuits and weights between them} are parameterised the same way --- a meta-organisation that might let us escape "the exponential trap of activation space" a second time, one level up. (ii) \textbf{Universality}: features and circuits (induction heads, curve detectors) recur across models, so understanding one model is a foothold in the next --- the difference between a "periodic table of features" and having to re-learn each model from scratch. (iii) The \textbf{micro-to-macro bridge}: induction heads are a microscopic phenomenon with a macroscopic signature (a visible bump in the loss curve, the 1-to-2-layer transition in scaling laws); perhaps every valley in the loss landscape is "a shadow cast by some circuit", and a dependency graph of circuits could explain learning dynamics and scaling laws in one theory. (iv) \textbf{Safety as a worst-case / out-of-distribution claim}: the "north star" is to certify things like "the model will never lie" --- which the essay explicitly notes is \emph{not} the expectation-over-a-distribution that ML normally handles, but a worst-case, out-of-distribution, possibly "deliberate-failure" claim; a \emph{complete} set of features might give a hook to enumerate over all model behaviour. The essay is candid that this last is "vague, highly speculative". It also flags the tension with \textbf{automated interpretability}: AI-automated audits might be the only way to scale, but invoking an AI to vouch for an AI is both aesthetically unsatisfying (cf. the computer proof of the four-colour theorem) and dangerous, because it correlates interpretability's failure modes with those of the systems it is meant to check.


\subsubsection*{The cold shower: \texorpdfstring{\textcite{toy-model-of-mechanistic-unfaithfulness}}{A Toy Model of Mechanistic (Un)Faithfulness}}

This note (also explicitly preliminary --- "treat these like a colleague's lab-meeting thoughts") is the group's only real mathematics, and it punctures a load-bearing assumption. The framing distinction is exactly the \emph{Landscape} survey IV.2 fork. When an SAE \textbf{models representation}, "much can be forgiven": the features are merely a basis for the activation space, so as long as they are monosemantic and reconstruct the activations, "you won't be misled" --- a basis cannot lie about what a vector contains. But a \textbf{transcoder models computation} --- it reads a layer's input and predicts its output --- and here a model can be "monosemantic, and \dots\ low MSE on distribution, but achieve this through a different computational mechanism", one that "may generalize differently off distribution". Every CS student knows many algorithms compute the same function (the paper's example: sorting algorithms); matching the function does not pin down the algorithm.

The construction is deliberately minimal. Target the map $y=\operatorname{abs}(x)$ coordinate-wise, which a clean ReLU network implements exactly as
\[
y_i \;=\; \operatorname{ReLU}(x_i) + \operatorname{ReLU}(-x_i).
\]
A transcoder trained on this (with a $\tanh$-$\ell_1$ penalty as a sharper surrogate for $\ell_0$) recovers exactly these two features per coordinate. Now inject a \textbf{repeated datapoint}: a fixed vector $p=[1,1,1,0,0,\dots]$ appears a fraction \texttt{repeat\_frac} of the time. The transcoder then grows a spurious \textbf{datapoint feature}
\[
y \mathrel{+}= p\cdot\operatorname{ReLU}(\lambda\, p\!\cdot\! x - b),
\]
which fires \emph{on $p$} and \emph{produces $p$}. Datapoint features are legitimate when the underlying model genuinely memorised that point --- but here the target $\operatorname{abs}(\cdot)$ "has no special notion of $p$." The transcoder has memorised something the model never did: monosemantic, sparse, low-error, and \textbf{mechanistically unfaithful}. The paper identifies the two necessary conditions and shows it is their \emph{product} that governs onset: (1) a repeated datapoint in the training set, and (2) a sparsity penalty (memorising $p$ lets that input be handled more sparsely). The divergence is bounded --- because the replacement happens "one step at a time", each transcoder step must match the true residual stream after a single nonlinearity, which sharply limits how far the mechanism can wander (the sorting analogy: force two sorts to agree on memory state after \emph{every} step and they are likely the same algorithm) --- but, as constructed, the bound is not tight enough to prevent the failure.

The proposed partial fix is \textbf{Jacobian matching}: penalise the squared Frobenius distance between the layer's true Jacobian and the transcoder's,
\[
\bigl\lVert J_{\mathrm{MLP}} - J_{\mathrm{Transcoder}}\bigr\rVert_F^2,
\]
estimated cheaply (it is expensive in full) by its expectation over a random probe vector $v$, $\;\mathbb{E}_v\,K\,\lVert J_{\mathrm{MLP}}v - J_{\mathrm{Transcoder}}v\rVert_2^2$, computable with one backprop. The principle: the Jacobian, viewed as a function over datapoints, is a near-canonical fingerprint of the \emph{mechanism}. Raw weights are not canonical (permute neurons, inversely scale input/output weights) and need context about how the nonlinearity behaves; but a ReLU layer's Jacobian is a sum of input/output-weight outer products that also records \emph{when neurons co-activate}. Concretely, the faithful baseline has a diagonal Jacobian on $p$ (features independent), whereas the datapoint-feature model's Jacobian shows every feature's local activation "explained" by all the others --- a visible signature of the spurious mechanism. Adding the penalty removes the memorisation feature in the toy example. The paper is honest about the limits: the transcoder can "cheat" by \textbf{gradient masking} (large weights, very negative biases, so a feature has large gradient but barely activates and so barely pays the $\ell_1$ cost --- an exploit of the $\ell_0/\ell_1$ gap, mitigated by the $\tanh$ surrogate); and, more subtly, one may \emph{not want} to match the Jacobian perfectly, because \textbf{superposition "blurs" the Jacobian} --- a feature spread over ten neurons contributes one large singular value and nine small ones that belong to other features, so the small-singular-value tail is an artifact to be ignored, not modelled.

The conclusion refactors all SAE/transcoder worries into two axes: \textbf{mechanistic faithfulness} (is it the true mechanism? --- "cuts at the heart of everything we hope for") and \textbf{simplicity} (is the mechanism captured cleanly? --- merely "a tax on how hard it will be to understand"). Faithfulness, it stresses, "is probably a spectrum", not a binary.


\subsection*{Common mathematical core}

The shared content is modest but real, and it is about \emph{epistemic status} rather than technique. Three threads.

\textbf{(1) The epistemic status of a circuit/feature claim.} All three papers treat an interpretation as a \emph{claim that can be true or false}, not merely a useful summary. The vocabulary paper locates the unit of such a claim (an interpretable variable / feature direction); the dreams paper asks what makes a claim \emph{checkable} (reducibility to small circuits where it becomes "mathematical reasoning"); the unfaithfulness paper exhibits a claim that is empirically well-supported on-distribution yet false. Together they sharpen the realisation that the field's working validations --- monosemanticity, low reconstruction error, the IOI-style \emph{faithfulness / completeness / minimality} of \textcite{interpretability-in-the-wild-ioi} --- are all behavioural, and behaviour underdetermines mechanism.

\textbf{(2) Matching a function vs.\ matching a mechanism --- a Jacobian.} This is the core's one genuinely mathematical idea. Two maps can agree as functions on a data distribution and disagree as \emph{mechanisms}; the unfaithfulness paper's operationalisation is that mechanism is carried by the \textbf{derivative}, and matching derivatives is strictly stronger than matching values. The hierarchy is the natural one a mathematician would write down: match the function $f$ on-distribution $\Rightarrow$ match its Jacobian $Df$ (first-order local mechanism) $\Rightarrow$ match the full computational trajectory (every intermediate state). Each level constrains more and admits fewer impostor mechanisms; the "one step at a time" bound is precisely the observation that constraining the trajectory (residual stream after every nonlinearity) is what keeps the divergence small.

\textbf{(3) The axiomatic dream.} The unifying ambition is that circuit claims could be \emph{reduced} to mathematical reasoning --- circuits as the axioms/assembly to which higher claims appeal when challenged. This is a claim about the \emph{logical form} of the eventual theory (a foundation one falls back to for verification), and it is what makes "is this interpretation correct?" a question one could hope to settle rather than merely adjudicate by taste.


\subsection*{How they diverge}

The three are not three takes on one question; they occupy three different registers, and their disagreement is in \emph{altitude and mood}.

\begin{itemize}
\item \textbf{Foundational vocabulary} (\textcite{mech-interp-variables-bases}): backward-looking and definitional. It fixes the metaphor and the units (binary/memory/variable; privileged basis; polysemanticity-as-superposition) without asserting that the program will succeed. It is the most conservative of the three.
\item \textbf{Aspirational vision} (\textcite{interpretability-dreams}): forward-looking and optimistic, explicitly speculative. It assumes the hard problem (superposition) gets solved and reasons about the upside --- families, universality, the macro bridge, safety certificates. Its mood is "if all goes well, here is the cathedral."
\item \textbf{Cautionary counterexample} (\textcite{toy-model-of-mechanistic-unfaithfulness}): present-tense and deflationary. It does not assume the hard problem away; it shows that even \emph{after} you have clean monosemantic sparse features, you can still be wrong about the mechanism. Its mood is "here is a brick that does not fit."
\end{itemize}

The tension is productive: the dreams paper's safety north-star (enumerate a complete feature set to certify worst-case behaviour) is exactly the use-case the unfaithfulness paper most endangers, since an out-of-distribution safety claim is precisely where a mechanism that "generalises differently off-distribution" bites. One paper sells the dream of trustworthy enumeration; another shows the enumerated objects can be subtly false in the regime that matters most.


\subsection*{What they answer, and what they don't}

What they answer: they give the field a self-consistent \emph{vocabulary} and a clear \emph{statement of ambition}, and they isolate \emph{the} central failure mode of computation-modelling (mechanistic unfaithfulness) with a clean toy demonstration and a candidate partial fix.

What they leave open --- honestly, and by their own admission:
\begin{itemize}
\item \textbf{Faithfulness is largely unformalised.} The IOI-era \emph{faithfulness/completeness/minimality} are behavioural and per-circuit; \emph{mechanistic} faithfulness is acknowledged to be "a spectrum" with no metric on it. Jacobian distance is a candidate \emph{ingredient}, not a definition: the paper itself says one should not match it perfectly (superposition blur), which means the target is not "$J$-distance $=0$" but some unspecified discounted version. There is, as yet, no scalar that says \emph{how} mechanistically faithful a transcoder is.
\item \textbf{The dream rests entirely on resolving superposition.} Essentially every aspiration in the dreams paper is explicitly gated on first untangling superposition; until then the larger-scale structure, the universality catalogue, and especially the "complete set of features" needed for safety enumeration are conjectural. The cautionary paper shows that resolving superposition (clean monosemantic features) is \emph{necessary but not sufficient} --- you also need faithfulness.
\item \textbf{The automation/trust tension is unresolved.} If interpretability must be automated to scale, it risks correlating its failures with the very systems it audits; if it must stay human-checkable, it may not scale. The papers name this dilemma but do not dissolve it.
\item \textbf{"Variable" and "circuit" are still informal.} The vocabulary paper's central object --- an interpretable variable / feature --- has no agreed formal definition (the recurring corpus-wide admission of \emph{Landscape} survey Part V); the dreams paper's "circuit" is likewise a working notion, not a defined one.
\end{itemize}


\subsection*{Connections}

\textbf{To Group 1 (validation metrics).} This group's "mechanistic faithfulness" is the deeper cousin of the IOI circuit's \emph{faithfulness / completeness / minimality} (\textcite{interpretability-in-the-wild-ioi}). Those are behavioural: does the sub-circuit reproduce the model's \emph{outputs}? Mechanistic faithfulness asks whether it reproduces the model's \emph{algorithm}. The unfaithfulness paper's contribution is precisely to show a gap between the two --- a replacement that would pass a behavioural check yet fails mechanistically --- which is the \emph{Landscape} survey V "rigour gradient" and "proxy objectives" worries made concrete.

\textbf{To Groups 3 and 4 (the representation/computation fork).} The representation-vs-computation distinction here \emph{is} \emph{Landscape} survey IV.2. SAEs (Group 3) model representation and "can't really mislead you"; transcoders and attribution graphs (Group 4) model computation and can. The mathematical statement of the fork: representation-modelling constrains a single space (a basis for activations), whereas computation-modelling must match a \textbf{Jacobian}, a linear map \emph{between} spaces --- and matching a function is strictly weaker than matching its derivative. This group supplies the worked counterexample that the IV.2 fork is real and not bookkeeping.

\textbf{To the gauge/metric theme.} The vocabulary paper is the one that \emph{named} the \textbf{privileged basis} (\emph{Landscape} survey I.4): meaning would live on arbitrary directions were it not for the activation function singling out a basis. This group thus identifies the very substrate the rest of the corpus negotiates with --- and the gauge theme resurfaces inside the unfaithfulness paper's own argument, which explicitly notes that raw weights are \emph{not} canonical (permutation and inverse input/output scaling are gauge moves) and proposes the Jacobian as a more nearly canonical, gauge-robust fingerprint of the mechanism. That is the same instinct as the \emph{Landscape} survey III.2 search for gauge-invariant quantities, applied to mechanism rather than to feature geometry.


\subsection*{For the mathematician}

This is the group where rigour is most conspicuously \emph{absent} and therefore where it could most cleanly be injected. Three concrete openings.

\textbf{(1) A quantitative theory of mechanistic faithfulness.} The paper gives the qualitative hierarchy
\[
\underbrace{f \approx \tilde f \text{ on-distribution}}_{\text{function match}}
\;\prec\;
\underbrace{Df \approx D\tilde f}_{\text{Jacobian match}}
\;\prec\;
\underbrace{\text{all intermediate states agree}}_{\text{full-trajectory match}},
\]
but no error semantics. The mathematician's task: define a faithfulness functional --- presumably some data-weighted, superposition-discounted norm of $J_{\mathrm{MLP}} - J_{\mathrm{Transcoder}}$ (the paper's own caveat that the small-singular-value tail should be ignored is a hint that the right object is a \emph{spectrally truncated} or relative Jacobian distance, not the raw Frobenius norm) --- and then prove the controlling inequalities: given a bound on function-match error and on Jacobian-match error, bound the \emph{out-of-distribution} divergence of the two mechanisms (e.g.\ along straight-line or geodesic perturbations off the data manifold). That would turn "faithfulness is a spectrum" into a metric with error bars, and would make the "one step at a time" bound (trajectory-matching limits divergence) into a theorem rather than an intuition.

\textbf{(2) Formalising "feature" and "circuit".} The corpus uses both as primitives. A serviceable definition of \emph{feature} must survive the $GL(V)$ gauge of I.4 --- so it is probably not "a direction in $V$" simpliciter but a direction together with the chosen metric, or better a gauge-invariant object (a covector read against the model's own maps, in the spirit of the probe/virtual-weight invariants of \emph{Landscape} survey III). A definition of \emph{circuit} wants to be a subobject of the computational graph closed under the maps that actually carry signal (the virtual-weight composites $W_{\mathrm{in}}^B W_{\mathrm{out}}^A$ of \emph{ML~Concepts} \S1.3), with faithfulness a statement that restricting to it perturbs the relevant Jacobians by a controlled amount.

\textbf{(3) What an "axiom" of a circuit would actually be.} The dreams paper's slogan --- circuit statements as "questions of mathematical reasoning" --- invites a literal reading: in what formal system? The honest version is that a circuit claim is a statement of the form "on input region $R$, with attention frozen to $A_0$ (\emph{ML~Concepts} \S1.4) and norm denominators held constant, the composite of these specified low-rank maps equals the model's output map up to error $\varepsilon$." Spelling out the quantifiers (which inputs, which frozen quantities, which error norm) is the work; doing so would reveal that the "axiom" is really a \emph{quantified approximation guarantee}, not an equation --- and would expose exactly the three different "linearities" (exact-by-freezing, local-linearisation, learned-approximation) of \emph{Landscape} survey IV.3 that this group's faithfulness worry is implicitly about. A theory that assigned each kind of circuit claim its proper error semantics would be the mathematical core this group currently only dreams of.


\printbibliography
\end{document}
